\section{Signed diffuse data: strong truncations and the full renormalized problem}
\label{sec:r11-signed-diffuse-renormalized}

Keep the same original arbitrary $(n-1,n-1)$ background, affinoid
neighborhood and part form as above.  Let $\mu$ be any finite
signed energy-diffuse Radon measure on $F$, and put
$B=\|\mu\|_{\rm TV}$.  All currents, coefficients and tests are
original CLD objects.  The new argument uses the already proved
positive-data theorem for the \emph{datum} $|\mu|$, with the same
original background $S$.  It does not replace that background.

\begin{lemma}[Bounded energy tests and signed Radon functionals]
\label{lem:r11-bounded-H-tests-diffuse}
Every bounded $z\in H$ has original compact smooth core
approximations converging strongly in $H$, uniformly bounded
in absolute value, and along a subsequence quasi-everywhere.
If a finite signed Radon measure $\nu$ on $F$ has a continuous
core functional for the energy norm, then $|\nu|$ charges no
energy-polar set.  If $h\in H$ and $D(hS)=\nu$ on every original
compact smooth test, then
\[
       \mathcal E(h,z)=-\int\widetilde z\,d\nu
                     \qquad(z\in H\cap L^\infty),
       \qquad h=u_\nu .
 \tag{SD1}
\]
The last equality is the already defined transposition solution.
Neither conclusion assumes that $|\nu|$ has a finite-energy norm.
\end{lemma}

\noindent\textbf{Status: \texttt{NEW\_PROVED}} in the stated
original-core realization; the compact-core capacity formula is
an already specified \texttt{EXTERNAL\_KNOWN} form input.

\begin{proof}
For bounded $z$, start with any original core sequence $z_j\to z$
in $H$ and compose it with one fixed smooth normal clamp that
fixes a neighborhood of the range of $z$ and is bounded globally.
The gradient-module chain rule proves strong convergence of these
compositions to $z$.  They are still original compact smooth
functions because the clamp fixes zero.  A summable-capacity
subsequence converges quasi-everywhere.

Let $K$ be an energy-polar compact set and choose an ordinary
relatively open neighborhood $G$ with compact closure such that
$|\nu|(G\setminus K)<\epsilon$.
The regular compact-core capacity formula, coercivity, a smooth
scalar clamp, and a fixed original cutoff supported in $G$
give core functions $0\le a_j\le1$, equal to one near $K$,
supported in $G$, with $\mathcal E(a_j)\to0$.
Multiplication by that cutoff preserves the energy convergence
by (AE12).  Hence $\int a_j\,d\nu\to0$, while
\[
       \left|\nu(K)-\int a_j\,d\nu\right|
                            \le |\nu|(G\setminus K)<\epsilon.
\]
Let $\epsilon\downarrow0$.  The same argument applies to every
polar compact subset of the positive and negative Hahn sets.
Radon inner regularity of the two variations therefore proves
that $|\nu|$ vanishes on every polar Borel set.

For bounded $z\in H$ use the uniformly bounded core
approximation just constructed.  Energy continuity passes its
left side in the original identity; diffuseness and dominated
convergence pass the right side.  This proves (SD1)'s first
assertion.  For $f\in L^\infty(\lambda)$ take $z=Gf$, which
is bounded by the torsion estimate.  Symmetry gives
\[
 \int h f\,d\lambda=\mathcal E(h,Gf)
                         =-\int\widetilde{Gf}\,d\nu.
\]
Since $h\in L^2(\lambda)\subset L^1(\lambda)$, transposition
uniqueness gives $h=u_\nu$.
\end{proof}

The usual product rule extends to bounded members of $H$:
$d(z_1z_2)=z_1dz_2+z_2dz_1$.  Indeed the preceding uniformly bounded
core approximations and the original gradient module pass this
identity by dominated convergence against the fixed limiting
energy measures.  This fact will be used for a product of two
bounded energy functions, not only a smooth cutoff.

\begin{theorem}[Strong clipping convergence for signed nest data]
\label{thm:r11-signed-strong-truncations}
Let $\mu_j=\mathbf1_{F_j}\mu$ be the finite-energy nest
restrictions from (MD3), and put $u_j=-G\mu_j$.
For every $k>0$,
\[
                          T_ku_j\longrightarrow T_ku_\mu
                                      \quad\hbox{strongly in }H.
 \tag{SD2}
\]
The conclusion does not require monotonicity of the signed
potentials.
\end{theorem}

\noindent\textbf{Status: \texttt{NEW\_PROVED}}.

\begin{proof}
Let $w=-u_{|\mu|}=G|\mu|\ge0$ be the positive-data Green
potential, with its canonical representative.  Positivity and
order of the transposition map give
$|u_j|,|u_\mu|\le w$ quasi-everywhere for all $j$.
Put $a_M=|\mu|\{w>M\}$, so $a_M\to0$, and define
\[
          \zeta_M=\frac{T_{2M}w-T_Mw}{M}\in H,\qquad
          0\le\zeta_M\le1.
\]
It equals one on $\{w\ge2M\}$.  The exact positive-data
band identity and the closed clipping rule give
\[
 \mathcal E(\zeta_M)\le\frac{a_M}{M},\qquad
 d\zeta_M=\frac1M\mathbf1_{\{M<w<2M\}}\,dw.
 \tag{SD3}
\]
Here $dw$ on the printed bounded band means the compatible
clipped gradient $d(T_{2M+1}w)$ there; no globally Radon
first-order current of the unbounded $w$ is required.
Test the finite-energy equation for $u_j$ with
$\zeta_M^2T_ku_j\in H\cap L^\infty$.
The original closed gradient product and clipping rules give
\[
 \int\zeta_M^2\,d\Gamma_S(T_ku_j)
  =-\int\zeta_M^2 T_k\widetilde u_j\,d\mu_j
                  -2\int\zeta_M T_k\widetilde u_j\,
                                      d\Gamma_S(u_j,\zeta_M).
\]
The first term in absolute value is at most $k a_M$.
The cross energy is carried where $M<w<2M$, and there
$|u_j|<2M$, so its $u_j$ factor may be clipped at $2M$.
Cauchy--Schwarz and (MD5) therefore prove
\[
 \sup_j\Gamma_S(T_ku_j)\{w\ge2M\}
   \le k a_M+2k\sqrt{\mathcal E(\zeta_M)
                                 \sup_j\mathcal E(T_{2M}u_j)}
   \le k a_M+2k\sqrt{2B a_M}\longrightarrow0.
 \tag{SD4}
\]
This is the required uniform control of energy in the region
where positive and negative potentials could be large and cancel.

On $A_M=\{w\le M\}$ one has $|u_j-u_l|\le2M$.
The first-order clipping rule and (MD5), applied to the
finite-energy datum $\mu_j-\mu_l$, give
\[
 \|\mathbf1_{A_M}(du_j-du_l)\|^2
      \le\mathcal E(T_{2M+1}(u_j-u_l))
      \le(2M+1)\|\mu_j-\mu_l\|_{\rm TV}\longrightarrow0.
 \tag{SD5}
\]
Thus the localized gradients have a fixed limit in the
original Hilbert gradient module.  Its energy measure is
absolutely continuous with respect to $\lambda$.
For a smooth scalar $P$ with $P(0)=0$ and compactly supported
first derivative, bounded convergence against this fixed
limiting module measure proves that
$\mathbf1_{A_M}d(P(u_j))$ converges strongly.
Here $u_j\to u$ in $L^1(\lambda)$, and hence in measure;
continuity and boundedness of $P'$ supply the required
convergence of its scalar multipliers.
The energy of $P(u_j)$ outside $A_M$ is bounded by a fixed
multiple of the tail of a clipping $T_Lu_j$, where $L$
contains the support of $P'$.  Formula (SD4) makes it small
uniformly in $j$ as $M\to\infty$.
It follows that $P(u_j)$ is Cauchy in $H$; its $L^2$ limit
is $P(u)$, by boundedness and the $L^1$ convergence of $u_j$.

Finally approximate $T_k$ by smooth normal clamps
$P_\epsilon$ with uniform scalar error at most $\epsilon$,
whose derivatives differ from that of $T_k$ only where
$k-\epsilon<|s|<k+\epsilon$.  For $0<\epsilon<k/2$, the
finite-energy band formula and the source equation give
\[
 \begin{split}
 \mathcal E(T_ku_j-P_\epsilon(u_j))
  &\le\mathcal E(T_{k+\epsilon}u_j-T_{k-\epsilon}u_j)\\
  &=-\int(T_{k+\epsilon}\widetilde u_j
                 -T_{k-\epsilon}\widetilde u_j)\,d\mu_j
       \le2\epsilon B.
 \end{split}
 \tag{SD6}
\]
The smooth compositions are strongly convergent by the
preceding paragraph.  First let $j,l\to\infty$ and then
$\epsilon\downarrow0$ in the triangle inequality.
This makes $T_ku_j$ Cauchy in $H$.  Its already known
$L^2$ limit is $T_ku$, proving (SD2).
\end{proof}

\begin{theorem}[Full signed renormalized chain and energy distribution]
\label{thm:r11-signed-renormalized-chain}
For $u=u_\mu$, the increasing compatible energy measure
$\Gamma_S(u)=\lim_k\Gamma_S(T_ku)$ is sigma-finite, finite
on every bounded value strip, and zero on every exact scalar
level.  If $\Phi\in C^2(\mathbf R)$ has bounded first derivative
and compactly supported second derivative, then
\[
 D(\Phi(u)S)=\Phi'(\widetilde u)\mu
                           +\Phi''(\widetilde u)\Gamma_S(u)
 \tag{SD7}
\]
on all original compact smooth tests.  Every term is coefficient-
defined, and the right side is a finite signed Radon measure.
If in addition $\Phi'$ has compact support, then
$\Phi(u)-\Phi(0)\in H$.
For every $0\le a<b$,
\[
 \mathcal E(T_bu-T_au)
                   =-\int(T_b\widetilde u-T_a\widetilde u)\,d\mu.
 \tag{SD8}
\]
The scalar energy image has density
\[
 \widetilde u_*\Gamma_S(u)=k_\mu(t)\,dt,\qquad
 k_\mu(t)=
 \begin{cases}
       \mu\{\widetilde u<t\},&t<0,\\
       -\mu\{\widetilde u>t\},&t>0.
 \end{cases}
 \tag{SD9}
\]
This density is nonnegative and bounded by $B$ almost everywhere;
its value at zero is immaterial.  In particular all bounded-value
weighted energy measures used here are finite Radon measures.
\end{theorem}

\noindent\textbf{Status: \texttt{NEW\_PROVED}} for signed
finite diffuse data.  Unweighted energy need not be an ordinary
locally finite current, as the positive-data counterexample already
shows.

\begin{proof}
Compatibility and level zero follow from the $H$ clipping
identities exactly as in (RT7)--(RT9), without any sign
hypothesis on $u$.  Choose $M$ containing the support of
$\Phi''$ in its interior.  The finite-energy chain for $u_j$
reads
\[
 D(\Phi(u_j)S)=\Phi'(\widetilde u_j)\mu_j
                         +\Phi''(\widetilde u_j)\Gamma_S(T_Mu_j).
\]
The left side passes because $\Phi$ is Lipschitz and
$u_j\to u$ in $L^1(\lambda)$.  The source term passes in
variation by the already proved capacity convergence of $u_j$
to $u$, diffuseness of $|\mu|$, bounded continuity of $\Phi'$, and
total-variation convergence of $\mu_j$ to $\mu$.
Strong convergence in (SD2), the energy variation estimate
(KT2), and bounded convergence against the fixed limiting
energy measure pass the last term.  This proves (SD7).
For merely $C^2$ scalar functions, ordinary smoothing gives
the same limit, uniformly in the first two derivatives on
the bounded active range and uniformly in scalar values
up to the original Lipschitz error.
If $\Phi'$ is compactly supported, its normalized composition
is a smooth function of one bounded clipping and lies in $H$.

The same strong clipping convergence passes the finite-energy
band identity, proving (SD8).  Apply the argument separately
to the bounded one-sided bands
\[
 B^+_{a,b}(s)=\min((s-a)^+,b-a),\qquad
 B^-_{a,b}(s)=-B^+_{a,b}(-s).
\]
Their derivatives are the indicators of $(a,b)$ and
$(-b,-a)$, respectively.  Composing a strongly convergent
clipping with these finite hinge combinations is strongly
continuous: the closed gradient formula applies, and the
limiting $H$ energy is zero on their finitely many corner
levels by (RT1).  Hence
\[
 \begin{split}
 \Gamma_S(u)\{a<u<b\}&=-\int_a^b\mu\{u>t\}\,dt,\\
 \Gamma_S(u)\{-b<u<-a\}&=\int_a^b\mu\{u<-t\}\,dt.
 \end{split}
\]
There is no energy mass at zero or at the interval endpoints.
These formulas prove (SD9).  Nonnegativity follows because
it is the density of a positive measure; the bound follows
from total variation of $\mu$.
\end{proof}

\begin{proposition}[Finite global level measures and clipping curvatures]
\label{prop:r11-signed-global-curvature-truncation}
The Kato measures $L_t(u;S)$ are finite for every real $t$,
with
\[
 \begin{aligned}
 L_t(u;S)(F)&\le |\mu|\{u\le t\}&& (t<0),\\
 L_t(u;S)(F)&\le |\mu|\{u>t\}&& (t>0),\\
 L_0(u;S)(F)&\le B.
 \end{aligned}
 \tag{SD10}
\]
For $M>0$, the actual signed Radon curvature
$\nu_M=D(T_Mu\,S)$ satisfies
\[
 \begin{split}
 \nu_M&=\mathbf1_{\{-M<u\le M\}}\mu+L_{-M}(u;S)-L_M(u;S),\\
 \|\nu_M\|_{\rm TV}&\le B,\qquad
 \|\nu_M-\mu\|_{\rm TV}
                    \le2|\mu|\{|u|\ge M\}\longrightarrow0.
 \end{split}
 \tag{SD11}
\]
Both $\nu_M$ and every $L_t(u;S)$ are energy diffuse.
No finite-energy norm of $|\nu_M|$ is claimed.
\end{proposition}

\noindent\textbf{Status: \texttt{NEW\_PROVED}}.

\begin{proof}
Apply (SD7) to the smooth one-sided hinges.  Their positive
defects are
$\Phi_\epsilon''(u-t)\Gamma_S(u)$, supported in the scalar
strip $t<u<t+\epsilon$.  Formula (SD9), the nonnegative
unit-integral second derivative, and its bounded density
show that their masses are at most $B$.
For $t>0$ the sharper bound is $|\mu|\{u>t\}$; for
$t<0$ and sufficiently small $\epsilon$ it is
$|\mu|\{u<t+\epsilon\}$.
They converge on original compact tests to $L_t(u;S)$.
Positivity, compact smooth exhaustion and then
$\epsilon\downarrow0$ prove (SD10).

The scalar identity $T_Ms=h_{-M}(s)-h_M(s)$ gives the
first formula of (SD11).  The three source regions
$\{u\le-M\}$, $\{-M<u\le M\}$ and $\{u>M\}$
partition the finite-value locus.  Triangle inequalities
and (SD10) therefore give both variation estimates in
(SD11).  Since $T_Mu\in H$, its finite Radon curvature
is diffuse by Lemma~\ref{lem:r11-bounded-H-tests-diffuse}.

For any fixed $t$ and $M>|t|$, the exact scalar identity
\[
                    h_t(T_Mu)=h_t(u)-h_M(u)
 \tag{SD12}
\]
holds on the finite-value locus and hence in the coefficient-defined
$L^1$ sense.  Its left side belongs to $H$.  Its curvature is finite Radon
by (SD10) and the Kato decompositions, hence diffuse by
the same lemma.  Thus
\[
 L_t(u;S)-L_M(u;S)
       =D(h_t(T_Mu)S)-\mathbf1_{\{t<u\le M\}}\mu
\]
is a diffuse finite signed measure.  On any polar Borel
set, the mass of $L_t$ is therefore at most $L_M(F)$,
which tends to zero by (SD10).  This proves diffuseness
of $L_t$ without assuming it passes through an arbitrary
weak measure limit.
\end{proof}

\begin{lemma}[Finite signed Radon data in the $H$ level theorem]
\label{lem:r11-H-level-carrier-finite-signed}
Let $v\in H$ and suppose $D(vS)=\sigma$ is a finite signed
Radon measure on the original tests.  Then the conclusion
(RT2)--(RT3) holds, with a finite diffuse positive level
measure, without assuming finite energy integral of
$|\sigma|$.
\end{lemma}

\noindent\textbf{Status: \texttt{NEW\_PROVED}} as the required
signed-data strengthening of the earlier $H$ carrier lemma.

\begin{proof}
The datum is diffuse and the current equation extends to
all bounded $H$ tests by Lemma
\ref{lem:r11-bounded-H-tests-diffuse}.
Consequently the finite-energy scalar chain proof applies:
its test $\psi\Phi'(v)$ is bounded and lies in $H$.
The one-sided band tests are also bounded $H$ vectors and
give (SD9) with $v,\sigma$ in place of $u,\mu$.
Thus the smoothed hinge defects have uniformly bounded
total mass, and their positive limit $L_t(v;S)$ is finite.
The signed measure
$D(h_t(v)S)=\mathbf1_{\{v>t\}}\sigma+L_t(v;S)$ is finite
and comes from an $H$ vector.  It is therefore diffuse by
the preceding lemma, as is $L_t$.

The original current identity for this hinge source now
extends to every bounded $H$ test.  Take
$z=\psi g(v)$ with $\psi$ an original compact nonnegative
smooth function and $g$ smooth bounded, with bounded
derivative, nonnegative and zero near $t$.
Strong convergence of the smooth hinges in $H$ and bounded
convergence against the finite measure $|\sigma|$ pass
\[
 \int z\,dL_{t,\epsilon}
     =-\mathcal E(h_{t,\epsilon}(v),z)
              -\int z\Phi_\epsilon'(v-t)\,d\sigma
\]
to $\int z\,dL_t$.
Its left side is zero for small $\epsilon$ because
$g\Phi_\epsilon''=0$.  The same countable family of
nonnegative tests as in the proof of (RT3) shows that
$L_t$ is carried by $\{\widetilde v=t\}$.
\end{proof}

\begin{theorem}[Exact level carriers and coarea for all signed diffuse data]
\label{thm:r11-signed-diffuse-exact-level}
For every finite signed energy-diffuse datum and every real $t$,
\[
                    \mathbf1_{\{\widetilde u\ne t\}}L_t(u;S)=0.
 \tag{SD13}
\]
For every bounded compactly supported Borel scalar $\theta$,
\[
                \int\theta(t)L_t(u;S)\,dt
                         =\theta(\widetilde u)\Gamma_S(u)
 \tag{SD14}
\]
as finite signed Radon measures.  The hard-curvature
variation bounds (SD11) hold on the complete original
space and every actual moment refinement.
\end{theorem}

\noindent\textbf{Status: \texttt{NEW\_PROVED}}.

\begin{proof}
First identify $\nu_M=\mu$ on $A_M=\{|u|<M\}$.
Choose a smooth $\Phi$ with $\Phi'$ zero outside $[-M,M]$
and strictly positive on $(-M,M)$.
The identities $\Phi(T_Mu)=\Phi(u)$ and (RT9) make their
weighted energies agree.  Compare (SD7) with the bounded
$H$ chain for $T_Mu$ and its finite signed datum $\nu_M$.
The latter chain is valid by the bounded-test extension,
so $\Phi'(u)(\nu_M-\mu)=0$, proving the assertion on $A_M$.

Fix $t$, choose $M>|t|$, and apply the preceding $H$ level
lemma to $v=T_Mu$.  The level measure $L_t(v;S)$ is carried
by $\{v=t\}=\{u=t\}$.
Compare the curvature of the two sides of (SD12), and then
restrict their actual Radon measures to $A_M$.
The just proved identity $\nu_M=\mu$ on $A_M$ cancels the
source terms and gives
\[
     \mathbf1_{A_M}(L_t(u;S)-L_M(u;S))
                         =\mathbf1_{A_M}L_t(T_Mu;S).
\]
It follows that
$L_t(u;S)(A_M\cap\{u\ne t\})\le L_M(u;S)(F)\to0$.
The sets $A_M$ exhaust the finite-value locus, and $L_t$ is
diffuse, so it gives no mass to the omitted polar infinite-
value set.  This proves (SD13) by monotone measure convergence.

For smooth compactly supported $\theta$, take
$\Phi(s)=\int\theta(t)(s-t)^+dt$.
It has bounded first derivative and compactly supported
second derivative $\theta$, so (SD7) and Fubini prove
(SD14) on every original test.
The positive kernels $L_t$, their global bound by $B$,
and their Borel original-test pairings extend the equality
to nonnegative Borel $\theta$; subtract positive and negative
parts.  The right side is finite by (SD9).
No unbounded unweighted energy current or restricted
background closedness has been assumed.
\end{proof}

\begin{definition}[The signed renormalized zero-boundary class]
\label{def:r11-signed-renormalized-boundary}
For the specified finite signed diffuse datum, a candidate is
$v\in L^1(F,\lambda)$ with a finite energy-quasicontinuous
representative and $T_kv\in H$ for every $k>0$.
Use the compatible sigma-finite energy measure of its
clippings.  Require the original current identity
\[
 D(\Phi(v)S)=\Phi'(\widetilde v)\mu
                     +\Phi''(\widetilde v)\Gamma_S(v)
 \tag{SD15}
\]
for every $\Phi\in C^2(\mathbf R)$ with compactly supported
first derivative, and the radiation condition
$k^{-1}\mathcal E(T_{2k}v-T_kv)\to0$.
These conditions specify the boundary and solution class
without using the transposition identity in their definition.
\end{definition}

\begin{theorem}[The complete finite-diffuse relative renormalized Dirichlet problem]
\label{thm:r11-signed-renormalized-Dirichlet}
Every finite signed energy-diffuse Radon datum has exactly one
candidate in Definition~\ref{def:r11-signed-renormalized-boundary}.
It is the transposition solution $u_\mu$, and all its original
current, exact-level, coarea and global curvature-variation
conclusions above hold.  The construction and solution class
are independent of the chosen faithful reference measure.
\end{theorem}

\noindent\textbf{Status: \texttt{NEW\_PROVED}} for the
printed relative linear operator and boundary class, in
arbitrary ambient dimension and arbitrary original $S$ of
the specified degree.

\begin{proof}
Existence follows from (SD7), the truncation membership,
quasicontinuity, and the already proved signed radiation
estimate (KT11).  For uniqueness choose smooth odd scalar
clamps $\Phi_k$ fixing the interval $[-k,k]$, constant
outside $[-2k,2k]$, with $\Phi_k(0)=0$,
$0\le\Phi_k'\le1$ and
$|\Phi_k''|\le C/k$ on the two transition bands.
Then $|\Phi_k(s)|\le|s|$ and $\Phi_k(s)\to s$.
For a candidate $v$, its prescribed chain gives
\[
 \begin{split}
 \nu_k:=D(\Phi_k(v)S)
       &=\Phi_k'(v)\mu+\Phi_k''(v)\Gamma_S(v),\\
 \|\nu_k-\mu\|_{\rm TV}
       &\le |\mu|\{|v|>k\}
                +\frac Ck\mathcal E(T_{2k}v-T_kv)\to0.
 \end{split}
 \tag{SD16}
\]
The measures $\nu_k$ are finite signed Radon measures.
The potential $\Phi_k(v)=\Phi_k(T_{2k}v)$ belongs to $H$,
so Lemma~\ref{lem:r11-bounded-H-tests-diffuse} proves
both diffuseness of $\nu_k$ and
$\Phi_k(v)=u_{\nu_k}$.  This step does not infer a
finite-energy norm for $|\nu_k|$ from its signed core norm.
The established total-variation stability gives
$\|\Phi_k(v)-u_\mu\|_1\le2\|\nu_k-\mu\|_{\rm TV}\to0$.
Dominated convergence also gives $\Phi_k(v)\to v$ in
$L^1(\lambda)$.  Thus $v=u_\mu$ almost everywhere.
Equality of all clipped $H$ vectors identifies the canonical
representatives quasi-everywhere.  Reference independence
follows from the already proved independence of the original
core form and the transposition solution.
\end{proof}

\paragraph{The remaining general problem.}
This section closes the signed diffuse chain and level-carrier
gates left open at the preceding positive-data stage.  It does
not establish ambient classical psh representability, ordinary
pointwise boundary regularity, or an unrestricted diagonal nonlinear Monge--Ampere
Dirichlet theory.  General analytic PS/PL relation descent and
low-real-degree CTO remain \texttt{OPEN}.  The background in
all the proved relative statements remains arbitrary.
Unrestricted equality of energy and BT polar ideals is already
\texttt{FALSE} by frozen R6, and unrestricted ambient classical
solvability is already \texttt{FALSE} by frozen R10; the larger
relative solution class is consistent with both obstructions.
